% !TEX root = ../main.tex
% บทที่ 2: ใช้ร่วมกับ main.tex และ styles/custom.sty

\begin{center}
{\LARGE\bfseries บทที่ 2}\\[0.5em]
{\LARGE\bfseries จำนวนเชิงซ้อน ฟังก์ชันตัวแปรเชิงซ้อน และเฟสเซอร์}\\[0.25em]
{\large\bfseries Complex Numbers, Complex Variable Functions and Phasors}
\end{center}

\noindent\textbf{ระยะเวลาที่ใช้สอน:} 2 สัปดาห์\\
\textbf{ชั่วโมงเรียน:} 6 ชั่วโมงบรรยาย + 4 ชั่วโมงปฏิบัติ\\
\textbf{บทที่รองรับ CLO:} CLO 1\\
\textbf{รายวิชา:} 5571102 คณิตศาสตร์วิศวกรรมไฟฟ้า\\
\textbf{Electrical Engineering Mathematics}

\section{สาระสำคัญ}

จำนวนเชิงซ้อน (Complex Number) เป็นเครื่องมือทางคณิตศาสตร์ที่มีความสำคัญอย่างยิ่งในงานวิศวกรรมไฟฟ้า โดยเฉพาะการวิเคราะห์วงจรไฟฟ้ากระแสสลับ (Alternating Current Circuit: AC Circuit) เนื่องจากสัญญาณไฟฟ้ากระแสสลับมีทั้ง “ขนาด” และ “มุมเฟส” ซึ่งไม่สามารถอธิบายได้ครบถ้วนด้วยจำนวนจริงเพียงอย่างเดียว จำนวนเชิงซ้อนช่วยให้สามารถแทนแรงดัน กระแส อิมพีแดนซ์ และกำลังไฟฟ้าในรูปแบบที่คำนวณได้สะดวกและเป็นระบบ (Kreyszig, 2011; Alexander \& Sadiku, 2021)

บทนี้เริ่มจากนิยามของจำนวนเชิงซ้อน การดำเนินการพื้นฐาน การแปลงระหว่างรูปเชิงพีชคณิต รูปเชิงขั้ว และรูปเอ็กซ์โพเนนเชียล จากนั้นอธิบายสูตรของออยเลอร์ (Euler’s Formula) และฟังก์ชันตัวแปรเชิงซ้อนเบื้องต้น ก่อนเชื่อมโยงเข้าสู่แนวคิดของเฟสเซอร์ (Phasor) ซึ่งเป็นพื้นฐานสำคัญในการวิเคราะห์วงจรไฟฟ้ากระแสสลับ อิมพีแดนซ์ของตัวต้านทาน ตัวเหนี่ยวนำ และตัวเก็บประจุ รวมถึงการประยุกต์ใช้ในระบบไฟฟ้ากำลังและระบบควบคุม (Hayt et al., 2019; Nilsson \& Riedel, 2019)

\section{ผลลัพธ์การเรียนรู้ประจำบท}

เมื่อศึกษาบทนี้แล้ว นักศึกษาสามารถ

\begin{enumerate}
\item อธิบายความหมายของจำนวนเชิงซ้อน ส่วนจริง ส่วนจินตภาพ และหน่วยจินตภาพได้
\item ดำเนินการบวก ลบ คูณ หาร และหาคอนจูเกตของจำนวนเชิงซ้อนได้
\item แปลงจำนวนเชิงซ้อนระหว่างรูปเชิงพีชคณิต รูปเชิงขั้ว และรูปเอ็กซ์โพเนนเชียลได้
\item อธิบายและประยุกต์ใช้สูตรของออยเลอร์ในการแทนสัญญาณไซน์ได้
\item อธิบายความหมายของฟังก์ชันตัวแปรเชิงซ้อนเบื้องต้นได้
\item แปลงสัญญาณไซน์ในโดเมนเวลาให้อยู่ในรูปเฟสเซอร์ได้
\item คำนวณอิมพีแดนซ์ของ R, L และ C ในวงจรไฟฟ้ากระแสสลับได้
\item วิเคราะห์วงจร AC อย่างง่ายโดยใช้เฟสเซอร์และจำนวนเชิงซ้อนได้
\item ใช้โปรแกรมคำนวณทางคณิตศาสตร์ เช่น Python หรือ MATLAB/Octave เพื่อช่วยคำนวณและแสดงผลจำนวนเชิงซ้อนได้
\end{enumerate}

\section{คำสำคัญ}

\begin{longtable}{L{0.24\textwidth}L{0.28\textwidth}L{0.38\textwidth}}
\toprule
\textbf{คำไทย} & \textbf{คำอังกฤษ} & \textbf{ความหมายโดยย่อ} \\
\midrule
จำนวนเชิงซ้อน & Complex Number & จำนวนที่ประกอบด้วยส่วนจริงและส่วนจินตภาพ \\
ส่วนจริง & Real Part & ส่วนของจำนวนเชิงซ้อนที่อยู่บนแกนจริง \\
ส่วนจินตภาพ & Imaginary Part & ส่วนของจำนวนเชิงซ้อนที่คูณกับ $j$ \\
หน่วยจินตภาพ & Imaginary Unit & $j=\sqrt{-1}$ ใช้ในงานวิศวกรรมไฟฟ้า \\
คอนจูเกต & Complex Conjugate & จำนวนเชิงซ้อนที่เปลี่ยนเครื่องหมายส่วนจินตภาพ \\
ขนาด & Magnitude / Modulus & ระยะจากจุดกำเนิดถึงจุดบนระนาบเชิงซ้อน \\
อาร์กิวเมนต์ & Argument & มุมของจำนวนเชิงซ้อนกับแกนจริงบวก \\
รูปเชิงขั้ว & Polar Form & การแทนจำนวนเชิงซ้อนด้วยขนาดและมุม \\
รูปเอ็กซ์โพเนนเชียล & Exponential Form & การแทนจำนวนเชิงซ้อนด้วย $re^{j\theta}$ \\
สูตรของออยเลอร์ & Euler’s Formula & ความสัมพันธ์ $e^{j\theta}=\cos\theta+j\sin\theta$ \\
เฟสเซอร์ & Phasor & จำนวนเชิงซ้อนที่แทนขนาดและเฟสของสัญญาณไซน์ \\
อิมพีแดนซ์ & Impedance & ความต้านทานรวมในวงจร AC มีหน่วยเป็นโอห์ม \\
\bottomrule
\end{longtable}

\section{บทนำ}

ในการวิเคราะห์วงจรไฟฟ้ากระแสสลับ สัญญาณแรงดันและกระแสมักอยู่ในรูปฟังก์ชันไซน์ เช่น

\[
v(t)=V_m\cos(\omega t+\phi)
\]

โดยที่  
$v(t)$ คือ แรงดันไฟฟ้า ณ เวลา $t$ มีหน่วยเป็นโวลต์ (V)  
$V_m$ คือ ค่าสูงสุดของแรงดัน มีหน่วยเป็นโวลต์ (V)  
$\omega$ คือ ความถี่เชิงมุม มีหน่วยเป็นเรเดียนต่อวินาที (rad/s)  
$t$ คือ เวลา มีหน่วยเป็นวินาที (s)  
$\phi$ คือ มุมเฟส มีหน่วยเป็นองศาหรือเรเดียน

หากวิเคราะห์สัญญาณดังกล่าวในโดเมนเวลาโดยตรง การบวก ลบ หรือเปรียบเทียบสัญญาณหลายตัวที่มีมุมเฟสต่างกันจะทำได้ค่อนข้างยุ่งยาก แต่เมื่อใช้จำนวนเชิงซ้อนและเฟสเซอร์ จะสามารถแทนสัญญาณไซน์ด้วยจำนวนเชิงซ้อนเพียงตัวเดียว ทำให้การคำนวณวงจรไฟฟ้ากระแสสลับทำได้ง่ายขึ้นอย่างมาก (Alexander \& Sadiku, 2021; Hayt et al., 2019)

ตัวอย่างเช่น แรงดัน

\[
v(t)=100\cos(377t+30^\circ)
\]

สามารถแทนด้วยเฟสเซอร์ได้เป็น

\[
\mathbf{V}=100\angle 30^\circ
\]

หรือหากใช้ค่า RMS จะเขียนเป็น

\[
\mathbf{V}=\frac{100}{\sqrt{2}}\angle 30^\circ
\]

การแทนด้วยเฟสเซอร์ช่วยให้นักศึกษาสามารถวิเคราะห์ความสัมพันธ์ระหว่างแรงดัน กระแส และอิมพีแดนซ์ในวงจร AC ได้ในลักษณะเดียวกับการใช้กฎของโอห์ม แต่เปลี่ยนจากจำนวนจริงเป็นจำนวนเชิงซ้อน (Nilsson \& Riedel, 2019)

\section{จำนวนเชิงซ้อนและการดำเนินการ}

\subsection{นิยามจำนวนเชิงซ้อน}

จำนวนเชิงซ้อน คือ จำนวนที่เขียนได้ในรูป

\[
z=a+jb
\]

โดยที่  
$z$ คือ จำนวนเชิงซ้อน  
$a$ คือ ส่วนจริง (Real Part) เขียนแทนด้วย $\operatorname{Re}(z)$  
$b$ คือ สัมประสิทธิ์ของส่วนจินตภาพ เขียนแทนด้วย $\operatorname{Im}(z)$  
$j$ คือ หน่วยจินตภาพ มีสมบัติ $j^2=-1$

ในทางคณิตศาสตร์ทั่วไปมักใช้ $i=\sqrt{-1}$ แต่ในงานวิศวกรรมไฟฟ้านิยมใช้ $j$ แทนหน่วยจินตภาพ เพื่อหลีกเลี่ยงความสับสนกับ $i$ ซึ่งมักใช้แทนกระแสไฟฟ้า (Electric Current) (Kreyszig, 2011; Alexander \& Sadiku, 2021)

จำนวนเชิงซ้อนสามารถแสดงบนระนาบเชิงซ้อน (Complex Plane) หรือแผนภาพอาร์กองด์ (Argand Diagram) โดยแกนนอนคือแกนจริง และแกนตั้งคือแกนจินตภาพ จุด $z=a+jb$ จะอยู่ที่พิกัด $(a,b)$

\subsection{การบวกและการลบจำนวนเชิงซ้อน}

กำหนดให้

\[
z_1=a_1+jb_1
\]

และ

\[
z_2=a_2+jb_2
\]

การบวกและลบจำนวนเชิงซ้อนทำได้โดยรวมส่วนจริงกับส่วนจริง และรวมส่วนจินตภาพกับส่วนจินตภาพ ดังนี้

\[
z_1+z_2=(a_1+a_2)+j(b_1+b_2)
\]

\[
z_1-z_2=(a_1-a_2)+j(b_1-b_2)
\]

การบวกและลบจำนวนเชิงซ้อนมีความสำคัญในการรวมสัญญาณไฟฟ้าหลายตัวที่มีเฟสต่างกัน เช่น การรวมแรงดันในวงจร AC หรือการรวมกระแสในจุดต่อวงจร (Alexander \& Sadiku, 2021)

\subsubsection*{ตัวอย่างที่ 2.1 การบวกและลบจำนวนเชิงซ้อน}

กำหนดให้

\[
z_1=4+j3,\quad z_2=2-j5
\]

จงหา $z_1+z_2$ และ $z_1-z_2$

\textbf{วิธีทำ}

\[
z_1+z_2=(4+2)+j[3+(-5)]
\]

\[
z_1+z_2=6-j2
\]

และ

\[
z_1-z_2=(4-2)+j[3-(-5)]
\]

\[
z_1-z_2=2+j8
\]

\textbf{ตอบ}  
$z_1+z_2=6-j2$  
$z_1-z_2=2+j8$

\subsection{การคูณจำนวนเชิงซ้อน}

การคูณจำนวนเชิงซ้อนในรูปเชิงพีชคณิตทำได้โดยกระจายพจน์เหมือนพีชคณิตทั่วไป และใช้สมบัติ $j^2=-1$

\[
z_1z_2=(a_1+jb_1)(a_2+jb_2)
\]

\[
z_1z_2=(a_1a_2-b_1b_2)+j(a_1b_2+a_2b_1)
\]

การคูณจำนวนเชิงซ้อนมักพบในการวิเคราะห์กำลังเชิงซ้อน อิมพีแดนซ์ และการแปลงรูปสัญญาณในวงจรไฟฟ้ากระแสสลับ (Hayt et al., 2019)

\subsubsection*{ตัวอย่างที่ 2.2 การคูณจำนวนเชิงซ้อน}

กำหนดให้

\[
z_1=3+j4,\quad z_2=2+j5
\]

จงหา $z_1z_2$

\textbf{วิธีทำ}

\[
z_1z_2=(3+j4)(2+j5)
\]

\[
=3(2)+3(j5)+j4(2)+j4(j5)
\]

\[
=6+j15+j8+j^2 20
\]

เนื่องจาก $j^2=-1$

\[
=6+j23-20
\]

\[
=-14+j23
\]

\textbf{ตอบ}

\[
z_1z_2=-14+j23
\]

\subsection{คอนจูเกตของจำนวนเชิงซ้อน}

คอนจูเกตของจำนวนเชิงซ้อน $z=a+jb$ เขียนแทนด้วย $z^*$ หรือ $\overline{z}$ นิยามเป็น

\[
z^*=a-jb
\]

ผลคูณของจำนวนเชิงซ้อนกับคอนจูเกตของตนเองจะได้จำนวนจริงเสมอ

\[
zz^*=(a+jb)(a-jb)=a^2+b^2
\]

คอนจูเกตมีบทบาทสำคัญในการหารจำนวนเชิงซ้อน การหาขนาดกำลังสอง และการคำนวณกำลังเชิงซ้อนในระบบไฟฟ้ากำลัง เช่น

\[
S=VI^*
\]

โดยที่ $S$ คือ กำลังเชิงซ้อน $V$ คือ เฟสเซอร์แรงดัน และ $I^*$ คือ คอนจูเกตของเฟสเซอร์กระแส (Grainger \& Stevenson, 1994; Alexander \& Sadiku, 2021)

\subsection{การหารจำนวนเชิงซ้อน}

การหารจำนวนเชิงซ้อนทำได้โดยคูณทั้งเศษและส่วนด้วยคอนจูเกตของตัวส่วน

\[
\frac{z_1}{z_2}=\frac{a_1+jb_1}{a_2+jb_2}
\]

คูณด้วยคอนจูเกตของตัวส่วน

\[
\frac{z_1}{z_2}=\frac{(a_1+jb_1)(a_2-jb_2)}{(a_2+jb_2)(a_2-jb_2)}
\]

\[
\frac{z_1}{z_2}=\frac{(a_1+jb_1)(a_2-jb_2)}{a_2^2+b_2^2}
\]

การหารจำนวนเชิงซ้อนใช้บ่อยในการหากระแสจากสมการเฟสเซอร์ของกฎโอห์ม

\[
\mathbf{I}=\frac{\mathbf{V}}{\mathbf{Z}}
\]

โดยที่  
$\mathbf{I}$ คือ เฟสเซอร์กระแส มีหน่วยเป็นแอมแปร์ (A)  
$\mathbf{V}$ คือ เฟสเซอร์แรงดัน มีหน่วยเป็นโวลต์ (V)  
$\mathbf{Z}$ คือ อิมพีแดนซ์ มีหน่วยเป็นโอห์ม ($\Omega$)

(Alexander \& Sadiku, 2021; Nilsson \& Riedel, 2019)

\section{รูปแบบของจำนวนเชิงซ้อน}

\subsection{รูปเชิงพีชคณิต}

รูปเชิงพีชคณิต หรือรูปสี่เหลี่ยม (Rectangular Form) คือรูปที่เขียนจำนวนเชิงซ้อนเป็น

\[
z=a+jb
\]

รูปนี้เหมาะสำหรับการบวกและลบจำนวนเชิงซ้อน เพราะสามารถแยกส่วนจริงและส่วนจินตภาพได้โดยตรง (Kreyszig, 2011)

\subsection{ขนาดและมุมของจำนวนเชิงซ้อน}

จำนวนเชิงซ้อน $z=a+jb$ มีขนาดหรือโมดูลัสเป็น

\[
r=|z|=\sqrt{a^2+b^2}
\]

และมีมุมหรืออาร์กิวเมนต์เป็น

\[
\theta=\tan^{-1}\left(\frac{b}{a}\right)
\]

โดยที่  
$r$ คือ ขนาดของจำนวนเชิงซ้อน  
$\theta$ คือ มุมของจำนวนเชิงซ้อน วัดจากแกนจริงบวก  
$a$ คือ ส่วนจริง  
$b$ คือ ส่วนจินตภาพ

ในการคำนวณจริงควรพิจารณาจตุภาคของจุด $(a,b)$ ด้วย เพราะค่าของ $\tan^{-1}(b/a)$ อาจให้มุมที่ต้องปรับตามตำแหน่งของจำนวนเชิงซ้อนบนระนาบเชิงซ้อน (Kreyszig, 2011; Brown \& Churchill, 2014)

\subsection{รูปเชิงขั้ว}

จำนวนเชิงซ้อนสามารถเขียนในรูปเชิงขั้วได้เป็น

\[
z=r\angle \theta
\]

หรือ

\[
z=r(\cos\theta+j\sin\theta)
\]

โดยที่  
$r$ คือ ขนาดของจำนวนเชิงซ้อน  
$\theta$ คือ มุมเฟสหรืออาร์กิวเมนต์

รูปเชิงขั้วมีความเหมาะสมมากสำหรับการคูณ หาร และการวิเคราะห์เฟสเซอร์ เพราะขนาดและมุมเฟสสามารถมองเห็นได้ชัดเจน (Alexander \& Sadiku, 2021)

\subsection{รูปเอ็กซ์โพเนนเชียล}

จากสูตรของออยเลอร์

\[
e^{j\theta}=\cos\theta+j\sin\theta
\]

ดังนั้นจำนวนเชิงซ้อนสามารถเขียนเป็นรูปเอ็กซ์โพเนนเชียลได้ดังนี้

\[
z=re^{j\theta}
\]

รูปเอ็กซ์โพเนนเชียลมีความสำคัญในงานสัญญาณ ระบบควบคุม การวิเคราะห์ฟูเรียร์ และลาปลาซ เพราะสามารถเชื่อมโยงสัญญาณไซน์กับฟังก์ชันเอ็กซ์โพเนนเชียลเชิงซ้อนได้โดยตรง (Kreyszig, 2011; Oppenheim et al., 1997)

\subsection{ตารางสรุปการแปลงรูปจำนวนเชิงซ้อน}

\begin{longtable}{L{0.28\textwidth}L{0.30\textwidth}L{0.34\textwidth}}
\toprule
\textbf{รูปแบบ} & \textbf{สมการ} & \textbf{เหมาะสำหรับ} \\
\midrule
รูปเชิงพีชคณิต & $z=a+jb$ & การบวกและลบ \\
รูปเชิงขั้ว & $z=r\angle\theta$ & การอ่านขนาดและมุมเฟส \\
รูปตรีโกณมิติ & $z=r(\cos\theta+j\sin\theta)$ & การเชื่อมโยงกับเรขาคณิต \\
รูปเอ็กซ์โพเนนเชียล & $z=re^{j\theta}$ & การวิเคราะห์สัญญาณ ฟูเรียร์ และลาปลาซ \\
\bottomrule
\end{longtable}

\section{สูตรของออยเลอร์และความหมายทางวิศวกรรมไฟฟ้า}

สูตรของออยเลอร์เป็นความสัมพันธ์พื้นฐานระหว่างฟังก์ชันเอ็กซ์โพเนนเชียลเชิงซ้อนและฟังก์ชันตรีโกณมิติ

\[
e^{j\theta}=\cos\theta+j\sin\theta
\]

จากสมการนี้สามารถเขียนได้ว่า

\[
\cos\theta=\frac{e^{j\theta}+e^{-j\theta}}{2}
\]

และ

\[
\sin\theta=\frac{e^{j\theta}-e^{-j\theta}}{2j}
\]

สูตรของออยเลอร์มีความสำคัญอย่างมากในการวิเคราะห์สัญญาณไฟฟ้า เพราะสัญญาณไซน์และโคไซน์สามารถมองเป็นส่วนจริงหรือส่วนจินตภาพของสัญญาณเอ็กซ์โพเนนเชียลเชิงซ้อน เช่น

\[
v(t)=V_m\cos(\omega t+\phi)
\]

สามารถเขียนได้เป็น

\[
v(t)=\operatorname{Re}\left\{V_me^{j(\omega t+\phi)}\right\}
\]

โดยที่ $\operatorname{Re}\{\cdot\}$ หมายถึง การนำเฉพาะส่วนจริงของจำนวนเชิงซ้อน สมการนี้เป็นพื้นฐานของเฟสเซอร์และการวิเคราะห์สัญญาณไซน์ในวงจรไฟฟ้ากระแสสลับ (Oppenheim et al., 1997; Hayt et al., 2019)

\section{ฟังก์ชันตัวแปรเชิงซ้อนเบื้องต้น}

\subsection{ความหมายของฟังก์ชันตัวแปรเชิงซ้อน}

ฟังก์ชันตัวแปรเชิงซ้อน (Complex Variable Function) คือ ฟังก์ชันที่มีตัวแปรต้นเป็นจำนวนเชิงซ้อน เช่น

\[
w=f(z)
\]

โดยที่

\[
z=x+jy
\]

และ

\[
w=u+jv
\]

โดยที่  
$z$ คือ ตัวแปรเชิงซ้อนอินพุต  
$w$ คือ ผลลัพธ์เชิงซ้อน  
$x,y$ คือ ตัวแปรจริงของอินพุต  
$u,v$ คือ ส่วนจริงและส่วนจินตภาพของผลลัพธ์

ดังนั้นสามารถเขียนได้ว่า

\[
f(z)=u(x,y)+jv(x,y)
\]

แนวคิดของฟังก์ชันตัวแปรเชิงซ้อนมีความสำคัญในงานวิศวกรรมไฟฟ้าขั้นสูง เช่น การวิเคราะห์สนามแม่เหล็กไฟฟ้า การวิเคราะห์สัญญาณ การแปลงลาปลาซ และการวิเคราะห์ระบบควบคุมในโดเมน $s$ (Kreyszig, 2011; Brown \& Churchill, 2014)

\subsection{ตัวอย่างฟังก์ชันตัวแปรเชิงซ้อน}

ตัวอย่างฟังก์ชันตัวแปรเชิงซ้อน เช่น

\[
f(z)=z^2
\]

ถ้า

\[
z=x+jy
\]

จะได้ว่า

\[
f(z)=(x+jy)^2
\]

\[
=x^2+2jxy+j^2y^2
\]

เนื่องจาก $j^2=-1$

\[
f(z)=x^2-y^2+j2xy
\]

ดังนั้น

\[
u(x,y)=x^2-y^2
\]

และ

\[
v(x,y)=2xy
\]

ฟังก์ชันลักษณะนี้ช่วยให้นักศึกษาเข้าใจว่าฟังก์ชันเชิงซ้อนสามารถแยกเป็นส่วนจริงและส่วนจินตภาพได้ ซึ่งเป็นพื้นฐานของการศึกษาฟังก์ชันวิเคราะห์ (Analytic Function) ในระดับสูงขึ้น (Kreyszig, 2011)

\subsection{ความเชื่อมโยงกับโดเมน \texorpdfstring{$s$}{s} ในระบบควบคุมและวงจรไฟฟ้า}

ในหัวข้อการแปลงลาปลาซ ตัวแปร $s$ เป็นจำนวนเชิงซ้อนที่นิยามโดย

\[
s=\sigma+j\omega
\]

โดยที่  
$s$ คือ ตัวแปรเชิงซ้อนในโดเมนลาปลาซ  
$\sigma$ คือ ส่วนจริง เกี่ยวข้องกับการเพิ่มขึ้นหรือลดลงแบบเอ็กซ์โพเนนเชียล  
$\omega$ คือ ความถี่เชิงมุม เกี่ยวข้องกับการแกว่งของสัญญาณ  
$j$ คือ หน่วยจินตภาพ

ตัวแปร $s$ มีความสำคัญในการวิเคราะห์เสถียรภาพของระบบควบคุมและวงจรไฟฟ้า เช่น โพล (Pole) และซีโร (Zero) ของฟังก์ชันถ่ายโอนจะอยู่บนระนาบเชิงซ้อน ซึ่งส่งผลต่อพฤติกรรมของระบบ เช่น การสั่น การหน่วง และความเร็วในการตอบสนอง (Ogata, 2010; Nise, 2020)

\section{เฟสเซอร์และการแทนสัญญาณไซน์}

\subsection{แนวคิดของเฟสเซอร์}

เฟสเซอร์ (Phasor) คือ จำนวนเชิงซ้อนที่ใช้แทนสัญญาณไซน์ในสถานะคงตัว โดยตัดส่วนเวลา $e^{j\omega t}$ ออก เหลือเฉพาะขนาดและมุมเฟสของสัญญาณ ตัวอย่างเช่น

\[
v(t)=V_m\cos(\omega t+\phi)
\]

สามารถแทนด้วยเฟสเซอร์ได้เป็น

\[
\mathbf{V}=V_m\angle\phi
\]

ถ้าใช้ค่า RMS จะเขียนเป็น

\[
\mathbf{V}=V_{rms}\angle\phi
\]

โดยที่

\[
V_{rms}=\frac{V_m}{\sqrt{2}}
\]

ในงานวิเคราะห์วงจรไฟฟ้ากระแสสลับนิยมใช้ค่า RMS เพราะสัมพันธ์โดยตรงกับกำลังไฟฟ้าเฉลี่ยที่โหลดได้รับ (Hayt et al., 2019; Alexander \& Sadiku, 2021)

\subsection{การแปลงจากโดเมนเวลาเป็นเฟสเซอร์}

หากสัญญาณอยู่ในรูปโคไซน์

\[
v(t)=V_m\cos(\omega t+\phi)
\]

เฟสเซอร์ในรูปค่าสูงสุดคือ

\[
\mathbf{V}=V_m\angle\phi
\]

และเฟสเซอร์ในรูปค่า RMS คือ

\[
\mathbf{V}=V_{rms}\angle\phi
\]

ถ้าสัญญาณอยู่ในรูปไซน์ เช่น

\[
v(t)=V_m\sin(\omega t+\phi)
\]

สามารถแปลงเป็นโคไซน์ได้โดยใช้ความสัมพันธ์

\[
\sin\theta=\cos(\theta-90^\circ)
\]

ดังนั้น

\[
v(t)=V_m\cos(\omega t+\phi-90^\circ)
\]

เฟสเซอร์จึงเป็น

\[
\mathbf{V}=V_m\angle(\phi-90^\circ)
\]

การกำหนดมาตรฐานว่าจะใช้ไซน์หรือโคไซน์เป็นสัญญาณอ้างอิงต้องทำให้สอดคล้องกันตลอดทั้งการคำนวณ เพื่อหลีกเลี่ยงความผิดพลาดของมุมเฟส (Nilsson \& Riedel, 2019)

\subsection{ตารางสรุปการแทนสัญญาณด้วยเฟสเซอร์}

\begin{longtable}{L{0.37\textwidth}L{0.27\textwidth}L{0.27\textwidth}}
\toprule
\textbf{สัญญาณในโดเมนเวลา} & \textbf{เฟสเซอร์แบบค่าสูงสุด} & \textbf{เฟสเซอร์แบบ RMS} \\
\midrule
$v(t)=V_m\cos(\omega t+\phi)$ & $V_m\angle\phi$ & $\dfrac{V_m}{\sqrt{2}}\angle\phi$ \\
$i(t)=I_m\cos(\omega t+\theta)$ & $I_m\angle\theta$ & $\dfrac{I_m}{\sqrt{2}}\angle\theta$ \\
$v(t)=V_m\sin(\omega t+\phi)$ & $V_m\angle(\phi-90^\circ)$ & $\dfrac{V_m}{\sqrt{2}}\angle(\phi-90^\circ)$ \\
\bottomrule
\end{longtable}

\section{อิมพีแดนซ์ในวงจรไฟฟ้ากระแสสลับ}

\subsection{ความหมายของอิมพีแดนซ์}

อิมพีแดนซ์ (Impedance) คือปริมาณที่แสดงการต่อต้านการไหลของกระแสไฟฟ้าในวงจรไฟฟ้ากระแสสลับ เขียนแทนด้วย $Z$ และมีหน่วยเป็นโอห์ม $(\Omega)$

\[
Z=R+jX
\]

โดยที่  
$Z$ คือ อิมพีแดนซ์ มีหน่วยเป็นโอห์ม $(\Omega)$  
$R$ คือ ความต้านทาน (Resistance) มีหน่วยเป็นโอห์ม $(\Omega)$  
$X$ คือ รีแอกแตนซ์ (Reactance) มีหน่วยเป็นโอห์ม $(\Omega)$  
$jX$ คือ ส่วนจินตภาพของอิมพีแดนซ์

ถ้า $X>0$ วงจรมีลักษณะเหนี่ยวนำ  
ถ้า $X<0$ วงจรมีลักษณะเก็บประจุ

(Alexander \& Sadiku, 2021; Hayt et al., 2019)

\subsection{อิมพีแดนซ์ของตัวต้านทาน ตัวเหนี่ยวนำ และตัวเก็บประจุ}

\begin{longtable}{L{0.25\textwidth}L{0.35\textwidth}L{0.30\textwidth}}
\toprule
\textbf{อุปกรณ์} & \textbf{ความสัมพันธ์เวลา} & \textbf{อิมพีแดนซ์} \\
\midrule
ตัวต้านทาน $R$ & $v(t)=Ri(t)$ & $Z_R=R$ \\
ตัวเหนี่ยวนำ $L$ & $v(t)=L\dfrac{di(t)}{dt}$ & $Z_L=j\omega L$ \\
ตัวเก็บประจุ $C$ & $i(t)=C\dfrac{dv(t)}{dt}$ & $Z_C=\dfrac{1}{j\omega C}=-\dfrac{j}{\omega C}$ \\
\bottomrule
\end{longtable}

โดยที่  
$R$ คือ ความต้านทาน มีหน่วยโอห์ม $(\Omega)$  
$L$ คือ ความเหนี่ยวนำ มีหน่วยเฮนรี (H)  
$C$ คือ ความจุไฟฟ้า มีหน่วยฟารัด (F)  
$\omega$ คือ ความถี่เชิงมุม มีหน่วยเรเดียนต่อวินาที (rad/s)

อิมพีแดนซ์เหล่านี้ช่วยให้สามารถใช้กฎของโอห์มในรูปเฟสเซอร์ได้

\[
\mathbf{V}=\mathbf{I}\mathbf{Z}
\]

ซึ่งเป็นเครื่องมือหลักของการวิเคราะห์วงจร AC ในสถานะคงตัว (Nilsson \& Riedel, 2019)

\subsection{เฟสของแรงดันและกระแสใน R, L และ C}

\begin{longtable}{L{0.25\textwidth}L{0.35\textwidth}L{0.32\textwidth}}
\toprule
\textbf{อุปกรณ์} & \textbf{ความสัมพันธ์เฟส} & \textbf{ความหมาย} \\
\midrule
ตัวต้านทาน $R$ & แรงดันและกระแสเฟสเดียวกัน & ไม่มีการนำหน้าหรือล้าหลัง \\
ตัวเหนี่ยวนำ $L$ & แรงดันนำหน้ากระแส $90^\circ$ & กระแสล้าหลังแรงดัน \\
ตัวเก็บประจุ $C$ & กระแสนำหน้าแรงดัน $90^\circ$ & แรงดันล้าหลังกระแส \\
\bottomrule
\end{longtable}

ความเข้าใจเรื่องเฟสมีความสำคัญต่อการวิเคราะห์ค่าตัวประกอบกำลัง (Power Factor)
ระบบไฟฟ้ากำลัง และการแก้ไขเพาเวอร์แฟกเตอร์ในโรงงานอุตสาหกรรม
(Grainger \& Stevenson, 1994; Alexander \& Sadiku, 2021)

\section{ตัวอย่างโจทย์พร้อมวิธีทำ}

\subsection*{ตัวอย่างที่ 2.3 การแปลงจำนวนเชิงซ้อนจากรูปเชิงพีชคณิตเป็นรูปเชิงขั้ว}

\textbf{โจทย์}  
จงแปลงจำนวนเชิงซ้อน

\[
z=3+j4
\]

ให้อยู่ในรูปเชิงขั้ว

\textbf{วิธีทำ}

ขั้นที่ 1 หาขนาด

\[
r=\sqrt{a^2+b^2}
\]

\[
r=\sqrt{3^2+4^2}
\]

\[
r=\sqrt{9+16}=5
\]

ขั้นที่ 2 หามุม

\[
\theta=\tan^{-1}\left(\frac{b}{a}\right)
\]

\[
\theta=\tan^{-1}\left(\frac{4}{3}\right)
\]

\[
\theta\approx 53.13^\circ
\]

ดังนั้น

\[
z=5\angle 53.13^\circ
\]

\textbf{ตอบ}

\[
3+j4=5\angle 53.13^\circ
\]

\subsection*{ตัวอย่างที่ 2.4 การแปลงจำนวนเชิงซ้อนจากรูปเชิงขั้วเป็นรูปเชิงพีชคณิต}

\textbf{โจทย์}  
จงแปลง

\[
z=10\angle 30^\circ
\]

ให้อยู่ในรูป $a+jb$

\textbf{วิธีทำ}

ใช้สมการ

\[
z=r(\cos\theta+j\sin\theta)
\]

แทนค่า

\[
z=10(\cos30^\circ+j\sin30^\circ)
\]

โดยที่

\[
\cos30^\circ=0.866,\quad \sin30^\circ=0.5
\]

ดังนั้น

\[
z=10(0.866+j0.5)
\]

\[
z=8.66+j5
\]

\textbf{ตอบ}

\[
10\angle30^\circ=8.66+j5
\]

\subsection*{ตัวอย่างที่ 2.5 การแปลงสัญญาณไซน์เป็นเฟสเซอร์}

\textbf{โจทย์}  
กำหนดแรงดันไฟฟ้า

\[
v(t)=311\cos(314t+30^\circ)\ \text{V}
\]

จงเขียนเฟสเซอร์แรงดันในรูปค่าสูงสุดและรูปค่า RMS

\textbf{วิธีทำ}

ขั้นที่ 1 พิจารณารูปทั่วไป

\[
v(t)=V_m\cos(\omega t+\phi)
\]

เปรียบเทียบกับโจทย์ จะได้

\[
V_m=311\ \text{V}
\]

\[
\omega=314\ \text{rad/s}
\]

\[
\phi=30^\circ
\]

ขั้นที่ 2 เฟสเซอร์แบบค่าสูงสุด

\[
\mathbf{V}=311\angle30^\circ\ \text{V}
\]

ขั้นที่ 3 คำนวณค่า RMS

\[
V_{rms}=\frac{V_m}{\sqrt{2}}
\]

\[
V_{rms}=\frac{311}{\sqrt{2}}
\]

\[
V_{rms}\approx 220\ \text{V}
\]

ดังนั้นเฟสเซอร์แบบ RMS คือ

\[
\mathbf{V}=220\angle30^\circ\ \text{V}
\]

\textbf{ตอบ}  
เฟสเซอร์แบบค่าสูงสุดคือ $311\angle30^\circ\ \text{V}$  
เฟสเซอร์แบบ RMS คือ $220\angle30^\circ\ \text{V}$

\subsection*{ตัวอย่างที่ 2.6 การหาอิมพีแดนซ์ของวงจร R-L อนุกรม}

\textbf{โจทย์}  
วงจร R-L อนุกรมมี $R=20\ \Omega$, $L=0.1\ \text{H}$ ต่อกับแหล่งจ่ายไฟความถี่ $f=50\ \text{Hz}$ จงหาอิมพีแดนซ์ของวงจร

\textbf{วิธีทำ}

ขั้นที่ 1 หาความถี่เชิงมุม

\[
\omega=2\pi f
\]

\[
\omega=2\pi(50)
\]

\[
\omega=314.16\ \text{rad/s}
\]

ขั้นที่ 2 หารีแอกแตนซ์ของตัวเหนี่ยวนำ

\[
X_L=\omega L
\]

\[
X_L=314.16(0.1)
\]

\[
X_L=31.416\ \Omega
\]

ขั้นที่ 3 เขียนอิมพีแดนซ์รวม

\[
Z=R+jX_L
\]

\[
Z=20+j31.416\ \Omega
\]

ขั้นที่ 4 หาขนาดของอิมพีแดนซ์

\[
|Z|=\sqrt{20^2+31.416^2}
\]

\[
|Z|=\sqrt{400+986.96}
\]

\[
|Z|\approx 37.24\ \Omega
\]

ขั้นที่ 5 หามุมของอิมพีแดนซ์

\[
\theta=\tan^{-1}\left(\frac{31.416}{20}\right)
\]

\[
\theta\approx 57.52^\circ
\]

ดังนั้น

\[
Z=37.24\angle57.52^\circ\ \Omega
\]

\textbf{ตอบ}

\[
Z=20+j31.416\ \Omega
\]

หรือ

\[
Z=37.24\angle57.52^\circ\ \Omega
\]

\subsection*{ตัวอย่างที่ 2.7 การหากระแสในวงจร AC ด้วยเฟสเซอร์}

\textbf{โจทย์}  
แหล่งจ่ายแรงดันมีค่า

\[
\mathbf{V}=220\angle0^\circ\ \text{V}
\]

ต่อกับโหลดที่มีอิมพีแดนซ์

\[
Z=20+j31.416\ \Omega
\]

จงหากระแส $\mathbf{I}$

\textbf{วิธีทำ}

ขั้นที่ 1 แปลงอิมพีแดนซ์เป็นรูปเชิงขั้ว จากตัวอย่างก่อนหน้าได้ว่า

\[
Z=37.24\angle57.52^\circ\ \Omega
\]

ขั้นที่ 2 ใช้กฎของโอห์มในรูปเฟสเซอร์

\[
\mathbf{I}=\frac{\mathbf{V}}{Z}
\]

แทนค่า

\[
\mathbf{I}=\frac{220\angle0^\circ}{37.24\angle57.52^\circ}
\]

การหารจำนวนเชิงซ้อนในรูปเชิงขั้ว ให้หารขนาดและลบมุม

\[
\mathbf{I}=\frac{220}{37.24}\angle(0^\circ-57.52^\circ)
\]

\[
\mathbf{I}=5.91\angle -57.52^\circ\ \text{A}
\]

\textbf{ตอบ}

\[
\mathbf{I}=5.91\angle -57.52^\circ\ \text{A}
\]

หมายความว่า กระแสมีขนาด $5.91\ \text{A}$ และล้าหลังแรงดัน $57.52^\circ$ ซึ่งสอดคล้องกับวงจรที่มีลักษณะเหนี่ยวนำ

\section{การประยุกต์ใช้ในงานวิศวกรรมไฟฟ้า}

\subsection{การวิเคราะห์วงจรไฟฟ้ากระแสสลับ}

จำนวนเชิงซ้อนและเฟสเซอร์ช่วยให้การวิเคราะห์วงจร AC ทำได้ง่ายขึ้น โดยเปลี่ยนสมการเชิงอนุพันธ์ของตัวเหนี่ยวนำและตัวเก็บประจุให้เป็นสมการพีชคณิตในรูปอิมพีแดนซ์ เช่น

\[
Z_L=j\omega L
\]

และ

\[
Z_C=\frac{1}{j\omega C}
\]

ดังนั้นวงจร R, L และ C สามารถวิเคราะห์ได้ด้วยกฎของโอห์มและกฎของเคอร์ชอฟฟ์ในรูปเฟสเซอร์ (Hayt et al., 2019; Alexander \& Sadiku, 2021)

\subsection{การวิเคราะห์สัญญาณไฟฟ้าและรูปคลื่น}

สูตรของออยเลอร์ช่วยให้สามารถแทนสัญญาณไซน์ด้วยฟังก์ชันเอ็กซ์โพเนนเชียลเชิงซ้อน ซึ่งเป็นพื้นฐานของการวิเคราะห์ฟูเรียร์ การวิเคราะห์สเปกตรัม และการประมวลผลสัญญาณไฟฟ้า เช่น การแยกรูปคลื่นที่ไม่เป็นไซน์ออกเป็นองค์ประกอบฮาร์มอนิก (Oppenheim et al., 1997; Lathi, 2009)

\subsection{การวิเคราะห์ระบบไฟฟ้ากำลัง}

ในระบบไฟฟ้ากำลัง จำนวนเชิงซ้อนใช้ในการวิเคราะห์กำลังไฟฟ้าเชิงซ้อน

\[
S=P+jQ
\]

โดยที่  
$S$ คือ กำลังเชิงซ้อน มีหน่วยเป็นโวลต์-แอมแปร์ (VA)  
$P$ คือ กำลังจริง มีหน่วยเป็นวัตต์ (W)  
$Q$ คือ กำลังรีแอกทีฟ มีหน่วยเป็นวาร์ (var)

แนวคิดนี้ใช้ในการวิเคราะห์โหลด มุมเฟส ตัวประกอบกำลัง และการชดเชยกำลังรีแอกทีฟในระบบไฟฟ้ากำลัง (Grainger \& Stevenson, 1994; Stevenson, 1982)

\subsection{การวิเคราะห์ระบบควบคุม}

ในระบบควบคุม ตัวแปร $s=\sigma+j\omega$ ใช้อธิบายพฤติกรรมของระบบในโดเมนลาปลาซ ตำแหน่งโพลและซีโรบนระนาบเชิงซ้อนส่งผลต่อเสถียรภาพและการตอบสนองของระบบ เช่น ระบบที่มีโพลใกล้แกนจินตภาพอาจมีการสั่นหรือใช้เวลานานกว่าจะเข้าสู่สภาวะคงตัว (Ogata, 2010; Nise, 2020)

\subsection{การวิเคราะห์สนามไฟฟ้าและสนามแม่เหล็ก}

แม้จำนวนเชิงซ้อนจะใช้เด่นชัดในวงจร AC แต่ยังมีบทบาทในการแทนสนามแม่เหล็กไฟฟ้าที่เปลี่ยนแปลงตามเวลา โดยเฉพาะการวิเคราะห์คลื่นแม่เหล็กไฟฟ้าในรูปเฟสเซอร์ เช่น สนามไฟฟ้าและสนามแม่เหล็กที่มีการเปลี่ยนแปลงแบบไซน์ตามเวลา (Sadiku, 2018)

\section{ตารางสรุปสูตรสำคัญประจำบท}

\begin{longtable}{L{0.26\textwidth}L{0.37\textwidth}L{0.27\textwidth}}
\toprule
\textbf{หัวข้อ} & \textbf{สูตร} & \textbf{ความหมาย} \\
\midrule
จำนวนเชิงซ้อน & $z=a+jb$ & รูปเชิงพีชคณิต \\
คอนจูเกต & $z^*=a-jb$ & เปลี่ยนเครื่องหมายส่วนจินตภาพ \\
ขนาด & $|z|=\sqrt{a^2+b^2}$ & ระยะจากจุดกำเนิด \\
มุม & $\theta=\tan^{-1}(b/a)$ & มุมกับแกนจริงบวก \\
รูปเชิงขั้ว & $z=r\angle\theta$ & แสดงขนาดและมุม \\
รูปตรีโกณมิติ & $z=r(\cos\theta+j\sin\theta)$ & เชื่อมกับเรขาคณิต \\
สูตรออยเลอร์ & $e^{j\theta}=\cos\theta+j\sin\theta$ & เชื่อม exponential กับ sine/cosine \\
รูปเอ็กซ์โพเนนเชียล & $z=re^{j\theta}$ & ใช้ในสัญญาณและระบบ \\
เฟสเซอร์แรงดัน & $\mathbf{V}=V_{rms}\angle\phi$ & แทนแรงดัน AC \\
กฎโอห์มเฟสเซอร์ & $\mathbf{V}=\mathbf{I}Z$ & ใช้วิเคราะห์วงจร AC \\
อิมพีแดนซ์ R & $Z_R=R$ & ตัวต้านทาน \\
อิมพีแดนซ์ L & $Z_L=j\omega L$ & ตัวเหนี่ยวนำ \\
อิมพีแดนซ์ C & $Z_C=\dfrac{1}{j\omega C}$ & ตัวเก็บประจุ \\
กำลังเชิงซ้อน & $S=VI^*$ & วิเคราะห์กำลังในระบบ AC \\
\bottomrule
\end{longtable}

\section{กิจกรรมการเรียนรู้ในชั้นเรียนและกิจกรรมปฏิบัติการ}

\subsection{กิจกรรมที่ 1 แผนภาพ Argand Diagram}

\textbf{วัตถุประสงค์}  
ให้นักศึกษาเข้าใจความสัมพันธ์ระหว่างรูปเชิงพีชคณิตและตำแหน่งของจำนวนเชิงซ้อนบนระนาบเชิงซ้อน

\textbf{ขั้นตอนกิจกรรม}

\begin{enumerate}
\item แบ่งนักศึกษาเป็นกลุ่ม กลุ่มละ 3--4 คน
\item กำหนดจำนวนเชิงซ้อนให้แต่ละกลุ่ม เช่น $2+j3$, $-4+j2$, $-3-j5$, $5-j1$
\item ให้นักศึกษาวาดตำแหน่งของจำนวนเชิงซ้อนบนระนาบเชิงซ้อน
\item คำนวณขนาดและมุมของจำนวนเชิงซ้อนแต่ละตัว
\item เปรียบเทียบผลการคำนวณกับตำแหน่งบนกราฟ
\end{enumerate}

\textbf{ผลลัพธ์ที่คาดหวัง}  
นักศึกษาสามารถแปลงจำนวนเชิงซ้อนจากรูปเชิงพีชคณิตเป็นรูปเชิงขั้ว และอธิบายความหมายทางเรขาคณิตได้

\subsection{กิจกรรมที่ 2 การแปลงสัญญาณไซน์เป็นเฟสเซอร์}

\textbf{วัตถุประสงค์}  
ให้นักศึกษาเข้าใจการแทนสัญญาณ AC ด้วยเฟสเซอร์

\textbf{ขั้นตอนกิจกรรม}

\begin{enumerate}
\item ผู้สอนกำหนดสัญญาณแรงดันและกระแส เช่น
\[
v(t)=100\cos(377t+20^\circ)
\]
\[
i(t)=5\cos(377t-30^\circ)
\]
\item ให้นักศึกษาเขียนเฟสเซอร์ของแรงดันและกระแส
\item คำนวณความต่างเฟสระหว่างแรงดันและกระแส
\item อภิปรายว่าวงจรมีลักษณะเป็นโหลดแบบใด เช่น ตัวต้านทาน ตัวเหนี่ยวนำ หรือตัวเก็บประจุ
\end{enumerate}

\subsection{กิจกรรมที่ 3 ปฏิบัติการคำนวณจำนวนเชิงซ้อนด้วย Python}

\textbf{วัตถุประสงค์}  
ให้นักศึกษาใช้โปรแกรมคำนวณช่วยตรวจสอบผลลัพธ์ของจำนวนเชิงซ้อน

\textbf{ตัวอย่างโค้ด}

\begin{lstlisting}[style=pythonstyle]
import cmath
import math

z = 3 + 4j

r = abs(z)
theta_rad = cmath.phase(z)
theta_deg = math.degrees(theta_rad)

print("z =", z)
print("Magnitude =", r)
print("Angle =", theta_deg, "degrees")
print("Conjugate =", z.conjugate())
\end{lstlisting}

\textbf{คำถามหลังปฏิบัติการ}

\begin{enumerate}
\item ค่า Magnitude ที่ได้ตรงกับการคำนวณด้วยมือหรือไม่
\item มุมที่คำนวณได้อยู่ในจตุภาคที่ถูกต้องหรือไม่
\item Python ใช้สัญลักษณ์ $j$ ในการแทนจำนวนเชิงซ้อนอย่างไร
\end{enumerate}

อ้างอิงแนวคิดการใช้ Python เพื่อการคำนวณจากเอกสาร Python Documentation และ NumPy Documentation

\section{แบบฝึกหัดท้ายบท}

\subsection{ระดับพื้นฐาน}

\begin{enumerate}
\item จงหาส่วนจริงและส่วนจินตภาพของจำนวนเชิงซ้อนต่อไปนี้
  \begin{enumerate}
  \item $z=5+j2$
  \item $z=-3+j7$
  \item $z=4-j6$
  \end{enumerate}
\item กำหนดให้ $z_1=3+j4$ และ $z_2=2-j5$ จงหา
  \begin{enumerate}
  \item $z_1+z_2$
  \item $z_1-z_2$
  \item $z_1z_2$
  \item $z_1^*$
  \end{enumerate}
\item จงแปลงจำนวนเชิงซ้อนต่อไปนี้ให้อยู่ในรูปเชิงขั้ว
  \begin{enumerate}
  \item $z=1+j1$
  \item $z=3+j4$
  \item $z=-2+j2$
  \end{enumerate}
\item จงแปลงจำนวนเชิงซ้อนต่อไปนี้ให้อยู่ในรูปเชิงพีชคณิต
  \begin{enumerate}
  \item $z=10\angle0^\circ$
  \item $z=5\angle60^\circ$
  \item $z=8\angle-30^\circ$
  \end{enumerate}
\item จงเขียน $e^{j30^\circ}$ ให้อยู่ในรูป $a+jb$
\end{enumerate}

\subsection{ระดับประยุกต์}

\begin{enumerate}[resume]
\item กำหนดแรงดันไฟฟ้า
\[
v(t)=170\cos(377t+45^\circ)\ \text{V}
\]
จงเขียนเฟสเซอร์แรงดันในรูปค่าสูงสุดและค่า RMS
\item กำหนดกระแสไฟฟ้า
\[
i(t)=10\sin(314t+20^\circ)\ \text{A}
\]
จงเขียนเฟสเซอร์กระแสโดยใช้โคไซน์เป็นสัญญาณอ้างอิง
\item วงจร R-L อนุกรมมี $R=10\ \Omega$, $L=0.05\ \text{H}$, $f=50\ \text{Hz}$ จงหาอิมพีแดนซ์รวมของวงจร
\item วงจร R-C อนุกรมมี $R=100\ \Omega$, $C=100\ \mu\text{F}$, $f=50\ \text{Hz}$ จงหาอิมพีแดนซ์รวมของวงจร
\item แหล่งจ่ายแรงดัน $\mathbf{V}=220\angle0^\circ\ \text{V}$ ต่อกับโหลด $Z=30+j40\ \Omega$ จงหากระแส $\mathbf{I}$
\end{enumerate}

\subsection{ระดับวิเคราะห์}

\begin{enumerate}[resume]
\item อธิบายว่าเหตุใดการใช้เฟสเซอร์จึงช่วยให้การวิเคราะห์วงจรไฟฟ้ากระแสสลับง่ายกว่าการวิเคราะห์ในโดเมนเวลาโดยตรง
\item เปรียบเทียบลักษณะเฟสของแรงดันและกระแสในตัวต้านทาน ตัวเหนี่ยวนำ และตัวเก็บประจุ พร้อมยกตัวอย่างการประยุกต์ใช้งานจริง
\item ถ้าวงจรมีอิมพีแดนซ์ $Z=20-j15\ \Omega$ จงวิเคราะห์ว่าโหลดมีลักษณะเหนี่ยวนำหรือเก็บประจุ พร้อมอธิบายเหตุผล
\item จงวิเคราะห์ผลของการเพิ่มความถี่ต่ออิมพีแดนซ์ของตัวเหนี่ยวนำและตัวเก็บประจุ
\item ระบบควบคุมมีโพลที่ $s=-2+j5$ และ $s=-2-j5$ จงอธิบายความหมายเบื้องต้นของส่วนจริงและส่วนจินตภาพของโพลต่อพฤติกรรมของระบบ
\end{enumerate}

\section{คำถามทบทวนท้ายบท}

\begin{enumerate}
\item จำนวนเชิงซ้อนแตกต่างจากจำนวนจริงอย่างไร
\item เหตุใดงานวิศวกรรมไฟฟ้าจึงนิยมใช้ $j$ แทน $i$
\item รูปเชิงพีชคณิตและรูปเชิงขั้วเหมาะกับการคำนวณแบบใด
\item คอนจูเกตของจำนวนเชิงซ้อนมีประโยชน์อย่างไรในการหารจำนวนเชิงซ้อน
\item สูตรของออยเลอร์มีความสำคัญต่อการวิเคราะห์สัญญาณไฟฟ้าอย่างไร
\item เฟสเซอร์คืออะไร และใช้แทนสัญญาณชนิดใด
\item ค่า RMS มีความสำคัญอย่างไรในการวิเคราะห์วงจร AC
\item อิมพีแดนซ์แตกต่างจากความต้านทานอย่างไร
\item เหตุใดตัวเหนี่ยวนำจึงมีอิมพีแดนซ์เป็น $j\omega L$
\item เหตุใดตัวเก็บประจุจึงมีอิมพีแดนซ์เป็น $\dfrac{1}{j\omega C}$
\item วงจรที่มีอิมพีแดนซ์ส่วนจินตภาพเป็นบวกบอกลักษณะใดของวงจร
\item วงจรที่มีอิมพีแดนซ์ส่วนจินตภาพเป็นลบบอกลักษณะใดของวงจร
\item เฟสเซอร์ช่วยในการวิเคราะห์ระบบไฟฟ้ากำลังอย่างไร
\item ตัวแปร $s=\sigma+j\omega$ ในระบบควบคุมมีความหมายอย่างไร
\item โปรแกรมคำนวณทางคณิตศาสตร์ช่วยตรวจสอบคำตอบของบทนี้ได้อย่างไร
\end{enumerate}

\section{สรุปท้ายบท}

บทที่ 2 กล่าวถึงจำนวนเชิงซ้อน ฟังก์ชันตัวแปรเชิงซ้อน และเฟสเซอร์ ซึ่งเป็นพื้นฐานสำคัญของการวิเคราะห์วงจรไฟฟ้ากระแสสลับและระบบทางวิศวกรรมไฟฟ้า โดยสรุปได้ดังนี้

\begin{enumerate}
\item จำนวนเชิงซ้อนเขียนได้ในรูป $z=a+jb$ โดย $a$ คือส่วนจริง และ $b$ คือส่วนจินตภาพ
\item หน่วยจินตภาพในงานวิศวกรรมไฟฟ้าใช้ $j$ โดยมีสมบัติ $j^2=-1$
\item จำนวนเชิงซ้อนสามารถเขียนได้หลายรูปแบบ ได้แก่ รูปเชิงพีชคณิต รูปเชิงขั้ว รูปตรีโกณมิติ และรูปเอ็กซ์โพเนนเชียล
\item ขนาดของจำนวนเชิงซ้อนคือ $|z|=\sqrt{a^2+b^2}$ และมุมคือ $\theta=\tan^{-1}(b/a)$
\item สูตรของออยเลอร์ $e^{j\theta}=\cos\theta+j\sin\theta$ เป็นสะพานเชื่อมระหว่างจำนวนเชิงซ้อนกับสัญญาณไซน์
\item เฟสเซอร์ใช้แทนสัญญาณไซน์ในสถานะคงตัว โดยแสดงขนาดและมุมเฟสของสัญญาณ
\item อิมพีแดนซ์เป็นจำนวนเชิงซ้อนที่ใช้แทนการต้านทานในวงจร AC
\item อิมพีแดนซ์ของตัวต้านทานคือ $R$ ของตัวเหนี่ยวนำคือ $j\omega L$ และของตัวเก็บประจุคือ $\frac{1}{j\omega C}$
\item กฎของโอห์มในรูปเฟสเซอร์คือ $\mathbf{V}=\mathbf{I}Z$
\item จำนวนเชิงซ้อนยังเชื่อมโยงกับระบบควบคุมผ่านตัวแปร $s=\sigma+j\omega$ และเชื่อมโยงกับระบบไฟฟ้ากำลังผ่านกำลังเชิงซ้อน $S=P+jQ$
\end{enumerate}

\section{ข้อเสนอแนะรูปภาพ แผนภาพ และกราฟที่ควรมีในบทเรียน}

\begin{longtable}{L{0.08\textwidth}L{0.28\textwidth}L{0.29\textwidth}L{0.20\textwidth}}
\toprule
\textbf{ลำดับ} & \textbf{รูปภาพ/แผนภาพ/กราฟที่ควรมี} & \textbf{จุดประสงค์ทางการเรียนรู้} & \textbf{แหล่งที่มา/คำค้นที่แนะนำ} \\
\midrule
1 & แผนภาพระนาบเชิงซ้อน & แสดงแกนจริง แกนจินตภาพ และตำแหน่ง $z=a+jb$ & complex plane Argand diagram \\
2 & แผนภาพการแปลง $a+jb$ เป็น $r\angle\theta$ & อธิบายขนาดและมุมของจำนวนเชิงซ้อน & สร้างเองด้วย Python หรือ GeoGebra \\
3 & แผนภาพวงกลมหนึ่งหน่วยและสูตรออยเลอร์ & เชื่อมโยง $e^{j\theta}$, $\cos\theta$, $\sin\theta$ & Euler formula unit circle \\
4 & กราฟสัญญาณไซน์และเฟสเซอร์หมุน & อธิบายที่มาของเฟสเซอร์ & phasor representation of sinusoid \\
5 & แผนภาพเฟสแรงดันและกระแสใน R, L, C & เปรียบเทียบการนำหน้าและล้าหลังของเฟส & หนังสือวงจรไฟฟ้า หรือสร้างเองด้วย Python \\
6 & วงจร R-L และ R-C อนุกรมพร้อมเฟสเซอร์ไดอะแกรม & เชื่อมโยงทฤษฎีกับการวิเคราะห์วงจรจริง & RL circuit phasor diagram, RC circuit phasor diagram \\
7 & ระนาบ $s$ ในระบบควบคุม & แสดงตำแหน่งโพลและซีโรเบื้องต้น & s-plane poles zeros control system \\
8 & ตัวอย่างกราฟจำนวนเชิงซ้อนจาก Python & ฝึกใช้โปรแกรมคำนวณ & สร้างเองด้วย Python Matplotlib \\
\bottomrule
\end{longtable}

หมายเหตุ: รูปภาพสำหรับเอกสารประกอบการสอนควรสร้างเองจาก Python, MATLAB, GNU Octave หรือ GeoGebra เพื่อหลีกเลี่ยงปัญหาลิขสิทธิ์ และควรอ้างอิงแนวคิดจากตำราหลักที่ใช้ในบทเรียน

\section{เอกสารอ้างอิงท้ายบท}

\subsection*{เอกสารอ้างอิงหลัก}

Alexander, C. K., \& Sadiku, M. N. O. (2021). \textit{Fundamentals of Electric Circuits} (7th ed.). McGraw-Hill.

Brown, J. W., \& Churchill, R. V. (2014). \textit{Complex Variables and Applications} (9th ed.). McGraw-Hill.

Hayt, W. H., Kemmerly, J. E., \& Durbin, S. M. (2019). \textit{Engineering Circuit Analysis} (9th ed.). McGraw-Hill.

Kreyszig, E. (2011). \textit{Advanced Engineering Mathematics} (10th ed.). John Wiley \& Sons.

Nilsson, J. W., \& Riedel, S. A. (2019). \textit{Electric Circuits} (11th ed.). Pearson.

\subsection*{เอกสารอ้างอิงด้านสัญญาณ ระบบ และระบบควบคุม}

Lathi, B. P. (2009). \textit{Linear Systems and Signals} (2nd ed.). Oxford University Press.

Nise, N. S. (2020). \textit{Control Systems Engineering} (8th ed.). Wiley.

Ogata, K. (2010). \textit{Modern Control Engineering} (5th ed.). Prentice Hall.

Oppenheim, A. V., Willsky, A. S., \& Nawab, S. H. (1997). \textit{Signals and Systems} (2nd ed.). Prentice Hall.

\subsection*{เอกสารอ้างอิงด้านระบบไฟฟ้ากำลังและสนามแม่เหล็กไฟฟ้า}

Grainger, J. J., \& Stevenson, W. D. (1994). \textit{Power System Analysis}. McGraw-Hill.

Sadiku, M. N. O. (2018). \textit{Elements of Electromagnetics} (7th ed.). Oxford University Press.

Stevenson, W. D. (1982). \textit{Elements of Power System Analysis} (4th ed.). McGraw-Hill.

\subsection*{แหล่งเรียนรู้แบบเปิดและเครื่องมือคำนวณ}

All About Circuits. (n.d.). \textit{AC Circuit Analysis}. \url{https://www.allaboutcircuits.com/textbook/alternating-current/}

GNU Octave Project. (n.d.). \textit{GNU Octave Documentation}. \url{https://docs.octave.org/}

MIT OpenCourseWare. (n.d.). \textit{Circuits and Electronics}. Massachusetts Institute of Technology. \url{https://ocw.mit.edu/}

NumPy Developers. (n.d.). \textit{NumPy Documentation}. \url{https://numpy.org/doc/}

Python Software Foundation. (n.d.). \textit{Python Documentation}. \url{https://docs.python.org/3/}

SciPy Developers. (n.d.). \textit{SciPy Documentation}. \url{https://docs.scipy.org/doc/scipy/}

SymPy Development Team. (n.d.). \textit{SymPy Documentation}. \url{https://docs.sympy.org/}

\section{แนวทางการตรวจสอบความถูกต้องของเนื้อหา}

\begin{enumerate}
\item ตรวจสอบการใช้สัญลักษณ์ $j$ ให้สอดคล้องกับงานวิศวกรรมไฟฟ้า
\item ตรวจสอบหน่วยของตัวแปร เช่น V, A, $\Omega$, H, F, rad/s
\item ตรวจสอบการแปลงมุมระหว่างองศาและเรเดียน
\item ตรวจสอบว่าการใช้ค่า RMS และค่าสูงสุดไม่ปะปนกัน
\item ตรวจสอบเครื่องหมายของรีแอกแตนซ์ โดยตัวเหนี่ยวนำเป็น $+j\omega L$ และตัวเก็บประจุเป็น $-j/(\omega C)$
\item ตรวจสอบการแปลงไซน์เป็นโคไซน์ก่อนเขียนเฟสเซอร์
\item ตรวจสอบว่าการอ้างอิงในแต่ละหัวข้อสอดคล้องกับเนื้อหาที่กล่าวถึง
\end{enumerate}

\clearpage
